\documentclass[11pt]{article}

\usepackage[utf8]{inputenc}
\usepackage[T1]{fontenc}
\usepackage{lmodern}
\usepackage{geometry}
\usepackage{amsmath,amssymb,amsfonts,amsthm,mathtools}
\usepackage{bm}
\usepackage{enumitem}
\usepackage{microtype}
\usepackage{hyperref}
\usepackage{titlesec}
\usepackage{booktabs}
\usepackage{array}
\usepackage{setspace}
\usepackage{mathrsfs}
\usepackage{physics}

\geometry{margin=1in}
\hypersetup{colorlinks=true, linkcolor=black, urlcolor=blue, citecolor=black}
\setlist[enumerate]{topsep=0.4em,itemsep=0.35em}
\setlist[itemize]{topsep=0.3em,itemsep=0.25em}
\titleformat{\section}{\large\bfseries}{\thesection.}{0.5em}{}
\titleformat{\subsection}{\normalsize\bfseries}{\thesubsection.}{0.5em}{}
\setstretch{1.06}

\newcommand{\Aether}{\AE{}ther}
\newcommand{\Aflow}{\AE{}ther-flow}
\newcommand{\TheoryName}{\Aflow{} Interpretation of Relativity}
\newcommand{\TheoryProgram}{\Aether{} / \Aflow{} framework}
\newcommand{\ExplicitPackage}{\mathfrak P_{\mathrm{exp}}}
\newcommand{\PrimitiveFamily}{\mathfrak P_{\mathrm{dyn}}}
\newcommand{\StableBridgeObject}{\mathfrak B_{\mathrm{st}}}
\newcommand{\GateFive}{G_{5}}

\newtheorem{definition}{Definition}
\newtheorem{proposition}{Proposition}
\newtheorem{remark}{Remark}
\newtheorem{corollary}{Corollary}
\newtheorem{theorem}{Theorem}

\title{\textbf{The \TheoryName{}}\\[0.4em]
\large Bounded GR Derivation\\
\large Bridge-Object Synthesis}
\author{}
\date{}

\begin{document}

\maketitle

\begin{abstract}
The current routing layer of \TheoryName{} no longer licenses the old recursive instruction ``go deeper until GR appears.'' The benchmark exact-closure package is already the finished public front door, the same-package observer-side remainder burden is already settled on the active completed branch, and the live research question is now narrower: can the current primitive-reservoir line derive the adopted benchmark under the five benchmark derivation gates?

This note records the bounded synthesis needed before a derivation-attempt manuscript can answer that question honestly. It freezes the benchmark-facing bridge question, freezes the stable bridge object under test, lists the benchmark-facing gains already in hand, and states exactly which gates remain open. The bridge question is whether the current line contains an explicit substrate law that forces the primitive observer-reduction chain yielding the continuity-Hessian compressed core and the benchmark one-metric observer package. On the present routing layer, the stable bridge object under test is the explicit substrate equation / symmetry / constraint package together with the primitive reversible constrained dynamics family that forces it. The later exact-shell-affine-nucleus, completion-core, trace-core, selector-law, and shell-restriction results remain preserved as benchmark-facing compression history already in hand, but they do not become the default live object under test merely by being deeper.

The positive result of the synthesis is therefore bounded and operational. Gates (1)--(4) remain open as first-principles derivation burdens from explicit substrate structure, whereas Gate (5) is already settled at same-package observer-side scope and now functions as a monitor condition rather than the default bottleneck. The next primary manuscript after this note was accordingly one bounded derivation attempt, and the present synthesis should now be read as the frozen-object setup behind that attempt and its later no-go upgrade.
\end{abstract}

\section{Introduction}

The benchmark exact-closure package of \TheoryName{} is already fixed. It is the finished public theory statement of the repository, and GR is adopted there exactly as the operational relativistic sector with ordinary matter coupling. The post-bridge routing notes now make equally explicit that the current line has \emph{not} yet derived that benchmark from explicit substrate structure, and that another same-output deeper-origin relay is no longer the default next move.

The remaining question is no longer indefinite depth. It is a bounded benchmark question: whether the current primitive-reservoir line can derive the adopted benchmark under the five benchmark derivation gates. That change of question requires a change of manuscript form. Before writing a bounded derivation attempt, one must freeze the bridge object actually under test, separate live bridge-object support from preserved relay history, and state exactly which gates remain open.

\paragraph{Framework and claim boundary.}
\TheoryProgram{} denotes the broader \Aether{} / \Aflow{} framework built on the claims that the \Aether{} is the underlying four-dimensional substrate of reality, the \Aflow{} is its intrinsic ordered motion, observed three-dimensional space is the local experiential slice of that deeper substrate, S-time is the experienced order of change arising from matter, light, and the \Aflow{}, and observed expansion is the three-dimensional appearance of deeper four-dimensional ordered motion. Within that framework, \TheoryName{} denotes the exact-closure theory in which the effective gravitational dynamics are adopted to be exactly Einsteinian with universal matter coupling. In the active sequence, \emph{adoption} means use of that established relativistic dynamics without claiming substrate derivation, while \emph{derivation} is reserved for a first-principles recovery from explicit substrate variables.

The present note stays strictly inside that boundary. It does not claim that GR has already been derived from the deeper \Aether{} / \Aflow{} structure. Its narrower task is to synthesize the bounded program now required by the routing layer.

\section{The Bounded Benchmark Question}

\begin{definition}[Benchmark derivation question]
The \emph{bounded benchmark question} of the current primitive-reservoir line is:
\begin{quote}
Can the present line derive the adopted GR benchmark under the five benchmark derivation gates, without introducing a second operative metric, without enlarging the ordinary low-energy matter law, and without reopening the already settled same-package observer-side closure question as if it were still the default bottleneck?
\end{quote}
\end{definition}

\begin{definition}[Bridge question behind the bounded benchmark question]
The \emph{bridge question} feeding that bounded benchmark question is whether the present line supplies an explicit substrate law that forces the primitive observer-reduction chain yielding:
\begin{enumerate}
    \item the continuity-Hessian compressed core,
    \item the one operative metric and ordinary-matter route of the benchmark observer package,
    \item and the benchmark one-metric observer package as derivational output rather than preserved package data.
\end{enumerate}
\end{definition}

\begin{proposition}[Why the benchmark question is now bounded]
At current routing scope, another same-output deeper-origin relay that preserves the same continuity-Hessian compressed core and the same benchmark one-metric observer package does not by itself change benchmark-facing status.
\end{proposition}

\begin{proof}
The benchmark exact-closure package is already fixed as the public front door, the post-bridge note records that the same-package observer-side remainder gate is no longer the live bottleneck on the active completed branch, and the upstream compression theorem records that another deeper-origin relay preserving the same observer-relevant compressed core is benchmark-neutral at current scope. Therefore the next honest question is no longer whether the line can keep descending, but whether the current line can discharge the benchmark derivation gates on a bounded test.
\end{proof}

\section{The Stable Bridge Object Under Test}

\begin{definition}[Stable bridge object under test]
For the present bounded program, freeze the stable bridge object under test as
\[
\StableBridgeObject=(\ExplicitPackage,\PrimitiveFamily),
\]
where:
\begin{enumerate}
    \item \(\ExplicitPackage\) is the explicit substrate equation / symmetry / constraint package fixed by the explicit-package necessity / uniqueness theorem;
    \item \(\PrimitiveFamily\) is the primitive reversible constrained dynamics family that forces \(\ExplicitPackage\) on the active line.
\end{enumerate}
This pair is the stable bridge object under test for the bounded derivation program.
\end{definition}

\begin{remark}
The benchmark-facing reason for Definition 3 is not that the later relay chain never compressed further. It did. The point is narrower: on the current routing layer, the explicit substrate package together with its primitive-dynamics forcing family is the stable bridge object anchor that still governs bounded derivation routing. It is the first object whose role remains benchmark-stable under the reset, rather than merely being the deepest preserved relay name.
\end{remark}

\begin{remark}
The later exact-shell affine nucleus, completion-state shell-trace core, combined-response ambient package, selector-law split core, and shell-restriction layers remain preserved benchmark-facing compression history already in hand. They matter as evidence that the bridge-side compression path is real. They do \emph{not} become the default live object under test merely because they are deeper. Promotion from preserved compression history to the new default target would require a bounded gate-changing reason, not depth alone.
\end{remark}

\begin{proposition}[What counts as progress relative to the frozen bridge object]
After Definition 3, a new primary manuscript counts as progress on the bounded program only if it:
\begin{enumerate}
    \item derives the benchmark package from \(\StableBridgeObject\),
    \item proves that \(\StableBridgeObject\) cannot discharge one or more open benchmark gates on the current line,
    \item or proves a further bridge-object compression with immediate benchmark-facing consequence for that bounded verdict.
\end{enumerate}
Another same-output relay that merely renames a deeper explanatory layer while leaving \(\StableBridgeObject\) and the benchmark gate status unchanged is not a new bounded-program answer.
\end{proposition}

\section{Benchmark-Facing Gains Already in Hand}

The present bounded program does not begin from zero. It begins from a specific set of benchmark-facing gains already on record.

\begin{table}[ht]
\centering
\renewcommand{\arraystretch}{1.2}
\begin{tabular}{>{\raggedright\arraybackslash}p{0.24\textwidth} >{\raggedright\arraybackslash}p{0.43\textwidth} >{\raggedright\arraybackslash}p{0.25\textwidth}}
\toprule
\textbf{Cluster} & \textbf{Results already in hand} & \textbf{Bounded-program consequence} \\
\midrule
Benchmark gatekeeping layer & Derivation-gates note; post-bridge derivation-gate status note; upstream compression theorem; continuity-Hessian necessity / uniqueness theorem; continuity-Hessian core derivation theorem. & Fixes the benchmark standard, settles the same-package observer-side closure bottleneck, and compresses benchmark-facing status to the continuity-Hessian core. \\
Primitive-reservoir bridge-entry chain & Datum-origin theorem; observer-sector defect-fiber factorization theorem; one-metric / ordinary-matter output-map theorem; chain-forcing theorem. & Fixes the bridge-entry gains that produce the primitive observer-reduction chain used by the bounded derivation program. \\
Stable bridge-object cluster & Selector-existence / uniqueness theorem; stationary-substrate-construction theorem; explicit-package necessity / uniqueness theorem. & Fixes the stable bridge object under test and rules out treating the current bridge object as a merely ad hoc presentation. \\
Preserved deeper compression history & Primitive reversible constrained dynamics forcing theorem; exact-shell affine nucleus collapse theorem; direct exact-shell affine nucleus forcing theorem from minimal completion core; completion-core trace-collapse theorem; ambient trace-core forcing theorem; combined-response ambient forcing theorem; selector-law split-core collapse theorem; split-core forcing theorem from primitive substrate shell restriction. & Records a real bridge-side compression path already in hand, while leaving the bounded program answerable to the frozen bridge object and the five gates rather than to recursive depth alone. \\
\bottomrule
\end{tabular}
\caption{Benchmark-facing gains already in hand before the bounded derivation attempt.}
\label{tab:gains}
\end{table}

\begin{proposition}[Immediate consequence of the gains already in hand]
Table \ref{tab:gains} shows that the remaining open burden is no longer to recover benchmark-facing bridge-entry data or to prove that the line has any bridge-object structure at all. The remaining burden is whether the frozen bridge object \(\StableBridgeObject\) can actually discharge the still-open benchmark gates.
\end{proposition}

\begin{proof}
The benchmark gatekeeping layer already fixes the five-gate standard and closes the same-package observer-side remainder bottleneck at active-package scope. The bridge-entry chain already fixes the continuity-Hessian-to-observer-package route used by the bounded program. The stable bridge-object cluster then fixes the current object under test rather than leaving that object undefined. Therefore the next open question is no longer what the test object is, but whether that object derives the adopted benchmark.
\end{proof}

\section{Exactly Which Gates Remain Open}

\begin{definition}[Open, settled, and monitor status]
For the bounded program of this note:
\begin{enumerate}
    \item a gate is \emph{open} if no theorem yet derives its benchmark content from explicit substrate structure through \(\StableBridgeObject\);
    \item a gate is \emph{settled at same-package scope} if the active package already proves the required observer-side statement there;
    \item a gate is a \emph{monitor condition} if it is already settled at same-package scope but must remain satisfied by any later bounded derivation attempt.
\end{enumerate}
\end{definition}

\begin{table}[ht]
\centering
\renewcommand{\arraystretch}{1.2}
\begin{tabular}{>{\raggedright\arraybackslash}p{0.20\textwidth} >{\raggedright\arraybackslash}p{0.19\textwidth} >{\raggedright\arraybackslash}p{0.46\textwidth}}
\toprule
\textbf{Gate} & \textbf{Bounded-program status} & \textbf{Exact reason} \\
\midrule
Einsteinian observer dynamics & Open & The benchmark package and post-bridge same-package closure record the Einsteinian observer law at package scope, but no theorem yet derives that law from \(\StableBridgeObject\) as explicit substrate output rather than preserved benchmark data. \\
One operative metric with universal matter coupling & Open & The one operative metric and ordinary matter route are already preserved on the active line, but no theorem yet derives their benchmark status from \(\StableBridgeObject\) without importing them as fixed continuation data. \\
Ordinary GR gauge and two-tensor content & Open & The current package preserves the ordinary low-energy gauge and two-tensor content, but no theorem yet shows that \(\StableBridgeObject\) forces that benchmark content as first-principles output on the current line. \\
Exact null, proper-time, and redshift recovery & Open & The benchmark relativistic recovery package remains fixed through the same operative metric, but no theorem yet derives that recovery package from \(\StableBridgeObject\) as explicit substrate output. \\
No genuine observer-side remainder & Settled at same-package scope; monitor condition & The post-bridge continuity-Hessian closure result already closes the same-package observer-side remainder gate on the active completed branch. A bounded derivation attempt must preserve that closure and may not reintroduce a genuine observer-side remainder. \\
\bottomrule
\end{tabular}
\caption{Exact gate status after the bounded bridge-object synthesis.}
\label{tab:gates}
\end{table}

\begin{theorem}[Exact open-gate verdict after the synthesis]
After the present synthesis, Gates (1)--(4) remain open on the current line as first-principles derivation burdens from explicit substrate structure through \(\StableBridgeObject\). Gate \(\GateFive\) is no longer the default open bottleneck, because it is already settled at same-package observer-side scope and now survives only as a monitor condition on any later bounded attempt.
\end{theorem}

\begin{proof}
Table \ref{tab:gates} states the status gate by gate. For Gates (1)--(4), the benchmark package preserves the relevant observer-level content, but preservation is not yet derivation from explicit substrate structure through \(\StableBridgeObject\). Those gates therefore remain open on the bounded program. For Gate \(\GateFive\), the post-bridge continuity-Hessian closure theorem already settles the no-genuine-observer-side-remainder condition on the active completed branch. Hence \(\GateFive\) is no longer the default bottleneck, though any later derivation attempt must keep it closed.
\end{proof}

\begin{corollary}[Current scientific status]
The current line has not yet derived the adopted benchmark from explicit substrate structure. The correct present status is therefore neither \texttt{Derived} nor an indefinite frontier relay. It is a bounded program with a frozen bridge object, gains already in hand, Gates (1)--(4) still open, and Gate \(\GateFive\) already settled at same-package scope.
\end{corollary}

\section{The Next Admissible Verdict Manuscript}

\begin{definition}[Admissible bounded verdicts]
The next primary manuscript after the present synthesis should end in exactly one of the following verdicts.
\begin{enumerate}
    \item \texttt{Derived}: the manuscript proves that \(\StableBridgeObject\) discharges the open benchmark gates while preserving the already settled monitor condition on Gate \(\GateFive\).
    \item \texttt{Not Derived On Current Line}: the manuscript proves that the current line, as constrained by \(\StableBridgeObject\) and the recorded claim boundary, fails one or more open benchmark gates.
    \item \texttt{Suspended Pending New Law}: the manuscript proves that the present line has reached a point where an additional explicit substrate law is genuinely missing and no current derivation or no-go theorem follows without introducing that new law.
\end{enumerate}
\end{definition}

\begin{proposition}[Why these are the only honest next verdicts]
After the present synthesis, depth alone is not a verdict. The next primary manuscript must answer the bounded benchmark question, so the only honest endings are the three verdicts of Definition 7.
\end{proposition}

\begin{proof}
The bridge object is now frozen, the gains already in hand are already recorded, and the open gates are stated explicitly in Table \ref{tab:gates}. A later manuscript that merely adds depth without changing the bounded verdict would leave the benchmark question unanswered. Therefore the next primary manuscript must terminate in a derivation verdict, a no-derivation verdict on the current line, or a principled suspension verdict pending a genuinely new law.
\end{proof}

\section{Conclusion}

The positive result of this note is not a completed first-principles derivation of GR. It is the bounded synthesis the routing layer now required. The benchmark question has been frozen, the stable bridge object under test has been frozen as the explicit substrate package together with its primitive-dynamics forcing family, the benchmark-facing gains already in hand have been listed, and the open-gate status has been stated exactly.

That exact status is now clean. Gates (1)--(4) remain open as first-principles derivation burdens from explicit substrate structure through the frozen bridge object. Gate (5) is already settled at same-package observer-side scope and survives as a monitor condition rather than the default bottleneck. The old recursive-frontier instruction is therefore no longer the right description of the active research job.

The next honest burden was correspondingly narrow: write one bounded derivation attempt against the bridge object frozen here and let it end honestly in \texttt{Derived}, \texttt{Not Derived On Current Line}, or \texttt{Suspended Pending New Law}. That bounded attempt was then written, and the current-line obstruction / no-go note later upgraded its suspension verdict to the scoped verdict \texttt{Not Derived On Current Line}. The present manuscript should therefore now be read as the setup document for that frozen bridge object, not as the canonical downstream status page for the frozen line.

\end{document}
